\section{Matrix actions in tensor coordinates}

Let $\mathbb F$ be $\mathbb R$ or $\mathbb C$. For a finite set $X$, write
$\ell^2(X;\mathbb F)$ for its Euclidean space. For finite $X,Y$, let
$\Phi:\ell^2(X;\mathbb F)\otimes\ell^2(Y;\mathbb F)
    \longrightarrow\ell^2(X\times Y;\mathbb F)$ be the unitary linear isometry
which sends the tensor product of the standard basis vectors at $x,y$ to
the standard basis vector at $(x,y)$. These coordinate identities support
finite tensor decompositions as used
in~\cite[Proposition 5.7, lines 1311--1333]{Arad2016RigorousRG}.

\begin{theorem}[The action of a Kronecker matrix]
    \label{thm:euclidean_kronecker_intertwining}
    \lean{Matrix.toEuclideanCLM_kronecker_intertwining}
    \leanok
    If $A$ is a square matrix on $X$ and $B$ is a square matrix on $Y$,
    their Euclidean actions satisfy
    $\Phi\circ(A\otimes B)=(A\otimes_{\mathrm K}B)\circ\Phi$.
    Here the left tensor product is the tensor product of operators and
    $\otimes_{\mathrm K}$ denotes the Kronecker matrix.
\end{theorem}

\begin{proof}\leanok
    In the tensor product of the two standard orthonormal bases, the
    matrix of the tensor product of two linear maps is their Kronecker
    product. Applying the coordinate representation gives the claimed
    equality of continuous linear maps.
\end{proof}

\begin{theorem}[Relabeling a matrix action]
    \label{thm:euclidean_submatrix_intertwining}
    \lean{Matrix.toEuclideanCLM_submatrix_intertwining}
    \leanok
    Let $e:Z\longrightarrow V$ be a bijection between finite sets and let
    $P$ be a square matrix on $V$. Set $P^e_{zz'}=P_{e(z),e(z')}$ and
    $U:\ell^2(Z;\mathbb F)\longrightarrow\ell^2(V;\mathbb F)$ by
    $(Uf)(v)=f(e^{-1}(v))$. Then $U\circ P^e=P\circ U$.
\end{theorem}

\begin{proof}\leanok
    Relabel the summation index in matrix-vector multiplication by $e$.
\end{proof}

\begin{theorem}[A relabeled finite tensor sum]
    \label{thm:euclidean_reindexed_tensor_sum}
    \lean{Matrix.toEuclideanCLM_reindexed_kronecker_sum_intertwining}
    \leanok
    \uses{thm:euclidean_kronecker_intertwining, thm:euclidean_submatrix_intertwining}
    Let $I$ be finite, let $e:Z\longrightarrow X\times Y$ be a bijection,
    and let $A_i,B_i$ be square matrices on $X,Y$, respectively.
    Let $U$ be the relabeling isometry associated to $e$ and put
    $T=\Phi^{-1}\circ U$. Set $P=\sum_{i\in I}A_i\otimes_{\mathrm K}B_i$.
    Then $T\circ P^e=(\sum_{i\in I}A_i\otimes B_i)\circ T$.
    The index set and factor order are unchanged; finite coordinate
    sets and $I$ may be empty.
\end{theorem}

\begin{proof}\leanok
    Sum the Kronecker-action identity over $I$ and combine it with
    relabeling of the matrix action. Apply the injectivity of $\Phi$
    to obtain the equality through $T=\Phi^{-1}\circ U$.
\end{proof}

\begin{theorem}[Exchanging factors in an intertwining decomposition]
    \label{thm:tensor_intertwining_exchange}
    \lean{TensorProduct.intertwining_sum_mapL_commIsometry}
    \leanok
    Let $\mathcal E,\mathcal L,\mathcal R$ be real or complex inner product
    spaces, let $I$ be finite, and let
    $T:\mathcal E\longrightarrow\mathcal L\otimes\mathcal R$ be a unitary
    linear isometry. Suppose $K,A_i,B_i$ are continuous linear operators
    on $\mathcal E,\mathcal L,\mathcal R$, respectively, and
    $TK=(\sum_{i\in I}A_i\otimes B_i)T$.
    Let $S$ exchange the two tensor factors and put $T'=ST$. Then
    \begin{align}
        T'K          & =\Bigl(\sum_{i\in I}B_i\otimes A_i\Bigr)T', \\
        T'K(T')^{-1} & =\sum_{i\in I}B_i\otimes A_i.
    \end{align}
    No finite-dimensionality or completeness assumption is required.
\end{theorem}

\begin{proof}\leanok
    The tensor-factor exchange satisfies
    $S(A_i\otimes B_i)=(B_i\otimes A_i)S$ for every $i$.
    Sum this identity and apply the supplied intertwining relation.
    Evaluation at $(T')^{-1}y$ gives the conjugated identity.
\end{proof}

\begin{theorem}[A conjugated tensor sum with the middle factor first]
    \label{thm:middle_first_conjugated_tensor_sum}
    \lean{Matrix.toEuclideanCLM_middleFirst_conjugation_of_kronecker_sum}
    \leanok
    Let $Z,X,Y,I$ be finite sets and $e:Z\to X\times Y$ a bijection.
    Suppose a matrix $P$ on $Z$ satisfies
    \begin{align}
        P_{zz'}=\sum_{i\in I}A_i(e(z)_X,e(z')_X)B_i(e(z)_Y,e(z')_Y).
    \end{align}
    Let $T:\C^Z\to\C^Y\otimes\C^X$ send the standard vector at $z$ to
    $e_{e(z)_Y}\otimes e_{e(z)_X}$. Then
    \begin{align}
        TPT^{-1}=\sum_{i\in I}B_i\otimes A_i.
    \end{align}
    Empty coordinate and summation sets are allowed.
\end{theorem}
\begin{proof}\leanok
    \uses{thm:euclidean_reindexed_tensor_sum,thm:tensor_intertwining_exchange}
    First apply the reindexed tensor-sum identity with factor order $X,Y$.
    Compose its isometry with exchange of the two tensor factors and then
    with the inverse isometry on the right. No matrix factors commute in
    this argument; only their tensor positions are exchanged.
\end{proof}

\begin{lemma}[Recovering a whole matrix from a tensor sum]
    \label{lem:middle_first_tensor_sum_iff}
    \lean{Matrix.toEuclideanCLM_middleFirst_conjugation_eq_sum_iff}
    \leanok
    With the finite sets, bijection and canonical unitary $T$ of the
    preceding theorem, the identities
    \begin{align}
        TPT^{-1}=\sum_{i\in I}B_i\otimes A_i,
        \qquad
        P_{zz'}=\sum_{i\in I}A_i(e(z)_X,e(z')_X)B_i(e(z)_Y,e(z')_Y)
    \end{align}
    are equivalent. Empty coordinate and summation sets are allowed.
\end{lemma}
\begin{proof}\leanok
    \uses{thm:middle_first_conjugated_tensor_sum}
    One implication is the preceding conjugation identity. For the
    converse, conjugate by $T^{-1}$ and use the injectivity of the
    Euclidean matrix action.
\end{proof}

\begin{lemma}[The same global filter and target]
    \label{lem:global_filter_from_middle_first_decomposition}
    \lean{Matrix.global_filter_of_middleFirst_tensor_decomposition}
    \leanok
    Let $Z,X,Y,I$ be finite sets, let $e:Z\to X\times Y$ be a
    bijection, and let $T:\C^Z\to\C^Y\otimes\C^X$ be its canonical
    middle-first unitary. Supply matrices $W_i$ on $Y$, operators
    $R_i$ on $\C^X$, a subspace $G\subseteq\C^Z$, and a real number
    $\Delta$. Suppose
    \begin{align}
        K_0=\sum_{i\in I}W_i\otimes R_i
    \end{align}
    preserves $TG$ and $(TG)^\perp$, and satisfies
    \begin{align}
        \|v\|&\le\|K_0v\| &&(v\in TG),\notag\\
        \|K_0v\|^2&\le\Delta\|v\|^2 &&(v\in(TG)^\perp).
    \end{align}
    Define $K=T^{-1}K_0T$, let $M_i$ be the matrix of $R_i$, and let
    $P$ be the matrix of $K$. Then
    \begin{align}
        P_{zz'}=\sum_{i\in I}M_i(e(z)_X,e(z')_X)W_i(e(z)_Y,e(z')_Y).
    \end{align}
    Moreover $K$ preserves $G,G^\perp$ and satisfies the corresponding
    two inequalities on those subspaces. No second filter is chosen.
\end{lemma}
\begin{proof}\leanok
    \uses{lem:middle_first_tensor_sum_iff}
    The preceding equivalence gives the matrix identity. Unitary
    conjugation carries $TG$ back to $G$, carries its orthogonal
    complement back to $G^\perp$, and preserves all norms.
\end{proof}

\begin{theorem}[A Hermitian matrix on its orthonormal eigenbasis]
    \label{thm:euclidean_hermitian_eigenvector}
    \lean{Matrix.IsHermitian.toEuclideanCLM_eigenvectorBasis}
    \leanok
    Let $A$ be a Hermitian matrix on a finite set $X$. Write
    $(v_j)_{j\in X}$ for its chosen orthonormal eigenbasis and
    $\lambda_j\in\mathbb R$ for the corresponding eigenvalues. Its
    continuous action on $\ell^2(X;\mathbb F)$ satisfies
    $Av_j=\lambda_jv_j$.
\end{theorem}
\begin{proof}\leanok
    The finite-dimensional spectral theorem gives this matrix-vector
    identity. The Euclidean action has the same coordinates.
\end{proof}
